\section{Условие}

{\bfseries Цель работы:} приобретение практических навыков управления потоками
в операционной системе и обеспечения синхронизации между потоками.

{\bfseries Задание:} составить программу на языке Си, обрабатывающую данные в
многопоточном режиме. При обработке использовать стандартные средства создания
потоков операционной системы. Ограничение максимального количества потоков,
работающих в один момент времени, задаётся ключом запуска программы.
Количество потоков, используемое программой, демонстрируется стандартными
средствами операционной системы. Исследовать зависимость ускорения и
эффективности алгоритма от входных данных и количества потоков и объяснить
полученные результаты.

{\bfseries Вариант:} 10.

\begin{enumerate}
  \item Решить систему линейных уравнений $Ax = b$ методом Гаусса.
  \item Потоки создаются средствами POSIX (\texttt{pthread\_create}); их число
        задаётся обязательным ключом \texttt{-t}.
  \item Система читается со стандартного ввода (с клавиатуры или из файла через
        перенаправление) либо генерируется случайно (ключ \texttt{-g}) для замеров.
  \item Число потоков работающей программы показывается через \texttt{ps -L},
        \texttt{/proc/<pid>/status} и \texttt{top -H}.
\end{enumerate}

\begin{lstlisting}[language={},title={Распределение строк расширенной матрицы $[A\,|\,b]$ между потоками ($T = 3$)}]
строка   элементы                      владелец
  0      a00  a01  a02  ...  a0n | b0   поток 0
  1      a10  a11  a12  ...  a1n | b1   поток 1
  2      a20  a21  a22  ...  a2n | b2   поток 2
  3      a30  a31  a32  ...  a3n | b3   поток 0
  ...                                   поток i mod T

шаг k:  поток 0      выбирает ведущую строку и ставит её на место k
        барьер       ведущая строка готова
        каждый поток вычитает строку k из своих строк с номерами больше k
        барьер       столбец k исключён
\end{lstlisting}
